\section{为什么必须上拓扑、弱收敛与对偶}

\label{sec:0-2}

上一节已经把问题摆清楚了：当变量落在函数、测度或轨道空间里时，有限维直觉并没有错，只是不再完整。接下来的问题不是“要不要学更多数学”，而是“哪一类数学分别在补哪一类缺口”。如果不把这些工具和它们所解决的问题对应起来，后文的拓扑、测度、对偶和算子章节就会显得像并列堆料。

第一个缺口是极限。极小化序列、逼近解和离散化解首先都在谈“往哪里收敛”。在有限维里，闭球紧性、序列收敛和连续映射几乎总能接上；但在无穷维里，极限常常只能在弱拓扑或弱$^\ast$拓扑下获得，于是我们必须引入拓扑语言，特别是网、闭包和下半连续性。第一个问题因此是：一个极值点究竟以什么意义存在？

第二个缺口是最优性。只要未知量不再是有限维向量，“导数 = 梯度向量”的图景就不再可靠。泛函的导数属于对偶空间，约束的法向量会变成连续线性泛函，乘子则通过伴随算子与原变量相连。因此，对偶、支撑超平面、次微分和共轭理论并不是后加的抽象层，而是无穷维最优性条件的自然语言。第二个问题于是变成：局部最优与全局对偶之间，究竟由什么对象沟通？

第三个缺口是算法。有限维中熟悉的梯度法、投影法、ADMM 和 KKT 系统，在无穷维里依然有对应物，但它们的真正母语不再是矩阵，而是邻近算子、预解式、单调算子和原--对偶分裂。算法不是在理论之后才出现的附录，它本身就在逼迫我们采用算子观点。第三个问题于是是：怎样把“求最优解”统一改写为“求某个算子的零点或不动点”？

\subsection*{记号与环境约定}

全书默认在三层环境之间切换：
\[
\text{局部凸空间} \subset \text{自反 Banach 空间} \subset \text{Hilbert 空间}
\]
这里的包含不是集合论意义，而是“结构越来越强、可用工具越来越多”的层级。讨论分离与对偶时，局部凸空间往往已经足够；讨论极小值存在和弱紧性时，通常需要自反 Banach 空间；讨论投影、梯度、邻近算子和分裂算法时，则往往要回到 Hilbert 空间。后文凡涉及对偶配对，统一写作
\[
\langle x^\ast,x\rangle.
\]
若空间本身是 Hilbert 空间，则通过 Riesz 表示可以把对偶元素重新看成向量，但这种识别必须在明确声明环境后才允许使用。

\begin{remark}
你可以把这一节看成全书的类型系统。一个命题若只在 Hilbert 空间成立，就不能被写成一般 Banach 结论；一个极小值若只在弱拓扑下存在，就不能偷换成强收敛结论；一个乘子若属于对偶空间，就不能直接画成欧氏法向量。后文所有章节的技术细化，本质上都只是把这些环境约定展开。
\end{remark}

如果把第 \ref{sec:0-1} 节的四个失效点重新对齐，那么拓扑和弱收敛负责处理极限，对偶和次微分负责处理最优性，算子和分裂负责处理算法。于是整本书的结构就不再像若干独立模块的拼装，而像对一个中心问题的三次逼近：先问极限能否存在，再问最优性如何表达，最后问如何把它算出来。第 \ref{sec:0-3} 节接下来只做一件事：告诉不同读者该沿哪条线先走。
